%=============================================================================
% APPENDIX A: NOTATION AND CONVENTIONS
% DFD Unified Review
%=============================================================================

\section{Notation and Conventions}\label{app:notation}

This appendix provides a complete reference for all notation used in the review.
Consistent conventions facilitate reproducibility and comparison with other work.

%-----------------------------------------------------------------------------
% A.1 FUNDAMENTAL FIELDS AND PARAMETERS
%-----------------------------------------------------------------------------
\subsection{Fundamental Fields and Parameters}\label{app:fields}

\begin{table*}[t]
\centering
\caption{Primary field variables and coupling parameters in \DFD.}
\label{tab:fields}
\begin{tabular}{llll}
\toprule
\textbf{Symbol} & \textbf{Name} & \textbf{Definition/Value} & \textbf{Units} \\
\midrule
\multicolumn{4}{l}{\textit{Fundamental field}} \\
$\psi$ & Scalar refractive field & Primary gravitational d.o.f. & dimensionless \\
$n$ & Refractive index & $n = e^\psi$ & dimensionless \\
$\Phi$ & Effective potential & $\Phi = -c^2\psi/2$ & m$^2$/s$^2$ \\
\midrule
\multicolumn{4}{l}{\textit{Acceleration scales}} \\
$\astar$ & Characteristic gradient scale & $2a_0/c^2 \approx 2.7 \times 10^{-27}$ m$^{-1}$ & m$^{-1}$ \\
$a_0$ & MOND acceleration scale & $2\sqrt{\alpha}\,cH_0 \approx 1.2 \times 10^{-10}$ m/s$^2$ & m/s$^2$ \\
$\mathbf{a}$ & Physical acceleration & $\mathbf{a} = (c^2/2)\nabla\psi$ & m/s$^2$ \\
$a^2$ & Acceleration invariant & $a^2 \equiv \mathbf{a}\cdot\mathbf{a}$ & m$^2$/s$^4$ \\
\midrule
\multicolumn{4}{l}{\textit{Coupling constants}} \\
$\ka$ & Self-coupling parameter & $\ka = 3/(8\alpha) \approx 51.4$ & dimensionless \\
$\kalpha$ & Clock coupling & $\kalpha = \alpha^2/(2\pi) \approx 8.5\times 10^{-6}$ & dimensionless \\
$K_A$ & Effective clock coupling & channel-resolved; Eq.~\eqref{eq:K-general} & dimensionless \\
\midrule
\multicolumn{4}{l}{\textit{Interpolating function}} \\
$\mu(x)$ & Crossover function & $\mu \to 1$ ($x \gg 1$), $\mu \to x$ ($x \ll 1$) & dimensionless \\
$\nu(y)$ & Inverse function & $y = x\mu(x)$, $x = y\nu(y)$ & dimensionless \\
$x$ & Dimensionless argument & $x = |\nabla\psi|/\astar = a/a_0$ & dimensionless \\
\bottomrule
\end{tabular}
\end{table*}

%-----------------------------------------------------------------------------
% A.2 COORDINATE AND METRIC CONVENTIONS
%-----------------------------------------------------------------------------
\subsection{Coordinate and Metric Conventions}\label{app:coordinates}

\paragraph{Metric Signature.}
We use the $(-,+,+,+)$ (mostly positive) signature throughout:
\begin{equation}
ds^2 = -c^2\,dt^2 + dx^2 + dy^2 + dz^2 \quad \text{(Minkowski)}.
\end{equation}
This matches the convention of Misner, Thorne \& Wheeler~\cite{MTW} and is 
standard in gravitational physics.

\paragraph{Optical Metric.}
The optical line element takes the form:
\begin{equation}
d\tilde{s}^2 = -\frac{c^2\,dt^2}{n^2} + d\mathbf{x}^2, \qquad n = e^\psi.
\end{equation}
Light rays satisfy $d\tilde{s}^2 = 0$. The coordinate speed of light is 
$c/n = c\,e^{-\psi}$.

\paragraph{Spherical Coordinates.}
For spherically symmetric problems:
\begin{equation}
d\mathbf{x}^2 = dr^2 + r^2(d\theta^2 + \sin^2\theta\,d\phi^2).
\end{equation}
The radial acceleration magnitude is $a = (c^2/2)|d\psi/dr|$.

\paragraph{Index Conventions.}
\begin{itemize}
\item Greek indices $\mu, \nu, \ldots \in \{0,1,2,3\}$ for spacetime
\item Latin indices $i, j, \ldots \in \{1,2,3\}$ for spatial components
\item Repeated indices imply summation (Einstein convention)
\end{itemize}

%-----------------------------------------------------------------------------
% A.3 PHYSICAL CONSTANTS
%-----------------------------------------------------------------------------
\subsection{Physical Constants}\label{app:constants}

\begin{table}[h!]
\centering
\caption{Physical constants used in calculations. Values from CODATA 2018.}
\label{tab:constants-appendix}
\resizebox{\columnwidth}{!}{%
\begin{tabular}{llll}
\toprule
\textbf{Symbol} & \textbf{Name} & \textbf{Value} & \textbf{Units} \\
\midrule
$c$ & Speed of light & $2.99792458 \times 10^8$ & m/s \\
$G$ & Gravitational constant & $6.67430(15) \times 10^{-11}$ & m$^3$\,kg$^{-1}$\,s$^{-2}$ \\
$\hbar$ & Reduced Planck constant & $1.054571817 \times 10^{-34}$ & J\,s \\
$\alpha$ & Fine-structure constant & $7.2973525693(11) \times 10^{-3}$ & dimensionless \\
$\alpha^{-1}$ & Inverse $\alpha$ & $137.035999084(21)$ & dimensionless \\
$H_0$ & Hubble constant & $70 \pm 2$ & km\,s$^{-1}$\,Mpc$^{-1}$ \\
$M_\odot$ & Solar mass & $1.98841 \times 10^{30}$ & kg \\
$R_\odot$ & Solar radius & $6.9634 \times 10^8$ & m \\
$\mathrm{AU}$ & Astronomical unit & $1.495978707 \times 10^{11}$ & m \\
\bottomrule
\end{tabular}
}% end resizebox
\end{table}

\paragraph{Derived Quantities.}
\begin{align}
r_s &= \frac{2GM}{c^2} && \text{(Schwarzschild radius)} \\
\Phi_\odot/c^2 &= -\frac{GM_\odot}{c^2 r} && \text{(Solar potential)} \\
&\approx -9.87 \times 10^{-9} \text{ at 1 AU}
\end{align}

%-----------------------------------------------------------------------------
% A.4 PPN AND GRAVITATIONAL WAVE PARAMETERS
%-----------------------------------------------------------------------------
\subsection{Post-Newtonian and Gravitational Wave Parameters}\label{app:ppn}

\begin{table}[h!]
\centering
\caption{Post-Newtonian parameters. \DFD predictions match GR exactly.}
\label{tab:ppn}
\begin{tabular}{llll}
\toprule
\textbf{Parameter} & \textbf{Meaning} & \textbf{GR} & \textbf{DFD} \\
\midrule
$\gamma$ & Space curvature per unit mass & 1 & 1 \\
$\beta$ & Nonlinearity in superposition & 1 & 1 \\
$\xi$ & Preferred-location effects (PPN) & 0 & 0 \\
$\alpha_1$ & Preferred-frame (PFE) & 0 & 0 \\
$\alpha_2$ & PFE parameter 2 & 0 & 0 \\
$\alpha_3$ & PFE parameter 3 & 0 & 0 \\
$\zeta_1$--$\zeta_4$ & Violation of momentum conservation & 0 & 0 \\
\bottomrule
\end{tabular}
\end{table}

\paragraph{Gravitational Wave Parameters.}
DFD's GW sector is constructed as a \textit{minimal transverse-traceless sector} that reproduces GR exactly in the radiative zone. The scalar field $\psi$ affects source dynamics but not GW propagation (see Sec.~\ref{sec:gw-tt} for construction, Sec.~\ref{sec:gw-ct-proof} for rigorous proof):
\begin{itemize}
\item $c_T$: Tensor mode propagation speed. \DFD: $c_T = c$ exactly (by conformal structure).
\item $h_{+}, h_{\times}$: Plus and cross polarizations. \DFD: identical to GR (no scalar GW modes in far zone).
\item $\delta\hat{\varphi}_k$: ppE phase deformation at $k$-PN order. \DFD: $\delta\hat{\varphi}_k = 0$ for compact binary accelerations $\gg a_0$.
\end{itemize}

%-----------------------------------------------------------------------------
% A.5 CLOCK AND LPI PARAMETERS
%-----------------------------------------------------------------------------
\subsection{Clock and LPI Parameters}\label{app:clock_notation}

\begin{table}[h]
\centering
\caption{Clock comparison parameters and sensitivities.}
\label{tab:clock_params}
\resizebox{\columnwidth}{!}{%
\begin{tabular}{lll}
\toprule
\textbf{Symbol} & \textbf{Definition} & \textbf{Typical Value} \\
\midrule
$\xi_{\rm LPI}^{\rm res}$ & Residual cavity--atom LPI parameter & DFD: screened residual; GR: $0$ \\
$S_A^\alpha$ & $\alpha$-sensitivity of clock A & See Table~\ref{tab:sensitivities} \\
$K_A$ & Effective clock coupling & channel-resolved Eq.~\eqref{eq:K-general} \\
$\Delta K_{AB}$ & Differential coupling & $K_A - K_B$ \\
$y$ & Fractional frequency & $y = \Delta\nu/\nu$ \\
\bottomrule
\end{tabular}%
}
\end{table}

\begin{table}[h!]
\centering
\caption{$\alpha$-sensitivities for selected clock transitions.}
\label{tab:sensitivities}
\begin{tabular}{llll}
\toprule
\textbf{Clock} & \textbf{Transition} & $S^\alpha$ & \textbf{Reference} \\
\midrule
Cs hyperfine & $6S_{1/2}$ F=3$\to$4 & $+2.83$ & \cite{Flambaum2006} \\
Rb hyperfine & $5S_{1/2}$ F=1$\to$2 & $+2.34$ & \cite{Flambaum2006} \\
H maser & 1S hyperfine & $+2.00$ & \cite{Flambaum2006} \\
Sr optical & $^1S_0 \to {}^3P_0$ & $+0.06$ & \cite{Dzuba2018} \\
Yb$^+$ E2 & $^2S_{1/2} \to {}^2D_{3/2}$ & $+0.88$ & \cite{Dzuba2018} \\
Yb$^+$ E3 & $^2S_{1/2} \to {}^2F_{7/2}$ & $-5.95$ & \cite{Dzuba2018} \\
Al$^+$ & $^1S_0 \to {}^3P_0$ & $+0.008$ & \cite{Dzuba2018} \\
\bottomrule
\end{tabular}
\end{table}

%-----------------------------------------------------------------------------
% A.6 GALACTIC DYNAMICS NOTATION
%-----------------------------------------------------------------------------
\subsection{Galactic Dynamics Notation}\label{app:galactic_notation}

\begin{table}[h!]
\centering
\caption{Notation for galactic dynamics and rotation curves.}
\label{tab:galactic}
\begin{tabular}{lll}
\toprule
\textbf{Symbol} & \textbf{Definition} & \textbf{Units} \\
\midrule
$V_c$ & Circular velocity & km/s \\
$V_{\rm flat}$ & Asymptotic flat velocity & km/s \\
$V_{\rm bar}$ & Baryonic (Newtonian) velocity & km/s \\
$g_{\rm obs}$ & Observed centripetal acceleration & m/s$^2$ \\
$g_{\rm bar}$ & Baryonic gravitational acceleration & m/s$^2$ \\
$M_{\rm bar}$ & Total baryonic mass & $M_\odot$ \\
$\Sigma$ & Surface mass density & $M_\odot$/pc$^2$ \\
$\Upsilon_\star$ & Stellar mass-to-light ratio & $M_\odot/L_\odot$ \\
\bottomrule
\end{tabular}
\end{table}

\paragraph{Key Relations.}
\begin{align}
g_{\rm obs} &= \frac{V_c^2}{r} && \text{(centripetal acceleration)} \\
g_{\rm bar} &= \frac{GM_{\rm bar}(<r)}{r^2} && \text{(Newtonian gravity)} \\
V_{\rm flat}^4 &= GM_{\rm bar}\,a_0 && \text{(BTFR, deep-field limit)}
\end{align}

%-----------------------------------------------------------------------------
% A.7 UNIT CONVENTIONS
%-----------------------------------------------------------------------------
\subsection{Unit Conventions}\label{app:units}

\paragraph{SI Units.}
All equations in this review are written in SI units unless otherwise noted. 
This ensures dimensional transparency and direct comparison with experimental 
values.

\paragraph{Geometric Units.}
For some derivations, particularly those involving spacetime structure, 
it is convenient to set $G = c = 1$. In these ``geometric units'':
\begin{align}
[M] &= [L] = [T], \\
1\,M_\odot &= 1.477\,\text{km} = 4.926\,\mu\text{s}.
\end{align}
When geometric units are used, this is stated explicitly.

\paragraph{Natural Units.}
For quantum considerations, $\hbar = c = 1$ gives:
\begin{align}
[M] &= [L]^{-1} = [T]^{-1}, \\
1\,\text{eV} &= 5.068 \times 10^6\,\text{m}^{-1} = 1.519 \times 10^{15}\,\text{s}^{-1}.
\end{align}

\paragraph{Gaussian vs.\ SI Electromagnetism.}
For electromagnetic quantities, we use SI (rationalized) units. The 
fine-structure constant is:
\begin{equation}
\alpha = \frac{e^2}{4\pi\epsilon_0\hbar c} \approx \frac{1}{137}.
\end{equation}

%-----------------------------------------------------------------------------
% A.8 ABBREVIATIONS AND ACRONYMS
%-----------------------------------------------------------------------------
\subsection{Abbreviations and Acronyms}\label{app:acronyms}

\begin{table}[h!]
\centering
\caption{Frequently used abbreviations.}
\label{tab:acronyms}
\resizebox{\columnwidth}{!}{%
\begin{tabular}{ll}
\toprule
\textbf{Acronym} & \textbf{Meaning} \\
\midrule
\DFD & Density Field Dynamics \\
GR & General Relativity \\
PPN & Parametrized Post-Newtonian \\
LPI & Local Position Invariance \\
MOND & Modified Newtonian Dynamics \\
BTFR & Baryonic Tully-Fisher Relation \\
RAR & Radial Acceleration Relation \\
GW & Gravitational Wave \\
ppE & Parametrized Post-Einsteinian \\
EFT & Effective Field Theory \\
UV & Ultraviolet (high-energy) \\
CMB & Cosmic Microwave Background \\
BAO & Baryon Acoustic Oscillations \\
SPARC & Spitzer Photometry and Accurate Rotation Curves \\
LLR & Lunar Laser Ranging \\
VLBI & Very Long Baseline Interferometry \\
\bottomrule
\end{tabular}
}% end resizebox
\end{table}

%-----------------------------------------------------------------------------
% A.9 SIGN CONVENTIONS SUMMARY
%-----------------------------------------------------------------------------
\subsection{Sign Convention Summary}\label{app:signs}

For quick reference, the key sign conventions are:

\begin{tcolorbox}[colback=blue!5!white, colframe=blue!50!black, title=Sign Conventions]
\begin{itemize}
\item \textbf{Metric signature:} $(-,+,+,+)$
\item \textbf{Potential sign:} $\Phi < 0$ in gravitational wells
\item \textbf{Field sign:} $\psi > 0$ in gravitational wells (so $n > 1$)
\item \textbf{Relation:} $\Phi = -c^2\psi/2$, hence $\psi = -2\Phi/c^2 > 0$
\item \textbf{Acceleration direction:} $\mathbf{a} = -\nabla\Phi = (c^2/2)\nabla\psi$ points toward mass
\item \textbf{Curvature:} Not applicable (\DFD uses flat background)
\end{itemize}
\end{tcolorbox}

These conventions ensure consistency with both the Newtonian limit and 
standard GR formulations.
